\begin{center}
    {\large \textbf{Giorno 10} \textbar\ Turtle Point \textbar\ 18 Agosto 2024 \textbar\ \textbf{Gili Trawangan}, Lombok}
\end{center}

\noindent\rule{\textwidth}{0.4pt}
\begin{figure}[h!] % Portrait images below
    \centering
    % First portrait image (on the left)
    \begin{minipage}[h]{0.48\textwidth} % Set width equal to half the page
        \centering
        \includegraphics[width=\textwidth]{images/flags/Screenshot GiliT.png}
    \end{minipage}
    \hfill
    % Text
    \begin{minipage}[h]{0.48\textwidth}
        \raggedright
        \textbf{Superficie}: 4.567 km² \\
        \vspace{0.3cm}
        \textbf{Popolazione}: 3,3 milioni \\
        \vspace{0.3cm}
        \textit{L'isola}
    \end{minipage}
\end{figure}

\landscapeLargeMargin{images/day10/PXL_20240818_011730854.jpg}

\landscapeLargeMargin{images/day10/PXL_20240818_014114822.MP.jpg}

\landscapeLargeMargin{images/day10/BBC3898D-EFCF-4269-89FF-F7C0E8D60503_1_201_a.jpeg}

\landscapeLargeMargin{images/day10/9EDF5DB9-5C3D-4E1E-A735-C373730592D2_1_201_a.jpeg}

\landscapeLargeMargin{images/day10/PXL_20240818_024159516.jpg}

\landscapeLargeMargin{images/day10/PXL_20240818_030430140.jpg}

\landscapeLargeMargin{images/day10/PXL_20240818_032541690.jpg}

\landscapeLargeMargin{images/day10/PXL_20240818_035531663.MP.jpg}

\landscapeLargeMargin{images/day10/PXL_20240818_065151050.MP.jpg}

\landscapeLargeMargin{images/day10/PXL_20240818_075508360.MP.jpg}

\noindent 18 Agosto 2024 \newline

Seconda notte trascorsa al Villas Edenia e questa volta siamo entrambi svegliati dal Muezzin. Io per fortuna mi riaddormento subito — non so se per abitudine o per esasperazione — mentre Elisa fa più fatica a riprendere sonno. O forse si è solo lamentata più di me al mattino per farmelo credere.

La sera precedente avevamo invitato Enrico e Cristina a fare colazione con noi, dal momento che abbiamo due colazioni aggiuntive incluse. Brekkie breakfast per tutti, con uova strapazzate, spinaci, pomodori arrostiti, pancetta, funghi e pane tostato, tranne per Cri che prova la colazione spagnola a base di tortilla, bruschette con pomodoro, pancetta e salsa aioli. Ad accompagnare, succo d'arancia e Lombok coffee.

Non resisto alla tentazione e mi butto in piscina per sorseggiare comodamente l'ultima tazza di caffè sulla cesta galleggiante di bambù, seguito a ruota da Elisa. Su Instagram avevo trovato foto della stessa esperienza con la didascalia "il sogno di una vita." A me è bastato un caffè per togliermi la voglia. Di Instagram.

Prima di lasciare il complesso vado in reception a ringraziare per la colazione, poiché continuo a non dare per scontato il servizio in camera. Ne approfitto per comunicare che possono venire a ritirare i piatti ma che non c'è bisogno di rifare le stanze — informazione che tornerà utile più avanti. Chiedo infine di prenotare un calesse per l'indomani per andare a prendere il traghetto per Bali. Mi consigliano di prenotarlo alle 9:30 ma chiedo di anticipare di dieci minuti per sicurezza — altra informazione che tornerà utile più avanti.

Salutiamo e ci dirigiamo verso il lato est dell'isola, attraversando la strada principale che taglia Gili Trawangan in due. Arrivati alla costa proseguiamo verso nord e notiamo con piacere che ogni spiaggia espone il cartello "Turtle Point". Sono così frequenti da apparire quasi come uno scam: "sì sì, qui davanti si vedono le tartarughe, prendi un lettino, entra in acqua e nuota per otto chilometri." Apprezzo particolarmente lo stile dei centri balneari sulla spiaggia: sono caratteristici e ognuno diverso dall'altro, ricordano molto l'atmosfera trovata su Ilha Grande in Brasile, con strutture di legno e bambù, riferimenti costanti al chill \& relax e cocktail bar a non finire.

Percorriamo tutta la costa e ci lasciamo ispirare da un paio di posti: essendo in quattro dobbiamo accontentare le esigenze di tutti in termini di ombra, sole, musica e intrattenimento. Siamo però incredibilmente bravi a scegliere subito il posto adatto.

Mentre Enrico e Cristina vanno a noleggiare maschera e salvagente — oggi anche Cri è decisa a esplorare i fondali — io ed Eli ci accomodiamo sui lettini e, come mio solito, inizio a fraternizzare col ragazzo indonesiano al bar che gestisce la spiaggia. Mi dice di chiamarsi Luca, ma non mi lascia del tutto convinto. Dal momento che gli ho dato corda viene spesso a pulirci i lettini dalla sabbia, e qui decido di indagare più a fondo.

"I don't trust you, is your name really Luca?" 

"I like Luca." 

"Oh ok. So what's your real name?" 

"I'm Johan." 

"Nice to meet you, Johan."

Penso che questa abitudine di fornire ai turisti nomi più familiari e facili da pronunciare sia legata al fastidio di non essere mai ricordati e mai chiamati per nome. Ti dico che mi chiamo in un modo semplice, così non sei costretto a impegnarti, e io mi immedesimo in quel nome così da sentirmi riconosciuto. Forse è tutta una mia turba e quel "Luca" me l'ha detto per prendersi gioco di me, ma preferisco dedicare sempre un po' di impegno e provare a ricordarmi il nome delle persone che si stanno dando da fare per me.

Enrico e Cristina tornano con l'attrezzatura e oggi sono loro per primi a entrare in acqua. Io non riesco a stare fermo: un po' per il caldo, un po' perché fremo per immergermi, dal momento che questo dovrebbe essere lo spot migliore per vedere le tartarughe.

Mentre siedo sotto l'ombrellone noto che le mie braccia e le mie gambe sono diventate color rosso fuoco e inizio a prendermi male. Sono stato sotto il sole pochi minuti senza crema solare e sono già scottato? Come farò a fare snorkeling tutto il giorno? Improvvisamente realizzo che l'ombrellone è di tela rossa un po' consumata e fa passare i raggi solari creando una luce colorata sulla mia pelle. La giornata di snorkeling è salva, ma meglio mettere un altro strato di crema.

Enri e Cri tornano euforici: hanno visto una tartaruga gigante nelle acque basse! Enri dice addirittura che, insieme al Monte Bromo, questo sarà uno dei ricordi migliori della vacanza.

Io ed Elisa gli diamo subito il cambio ai lettini e praticamente ci lanciamo in acqua. Elisa prova la maschera a tutto viso di Cristina, ma l'esperimento non va per il meglio: il ricircolo d'aria funziona molto male e respira praticamente solo anidride carbonica. Dopo dieci minuti di fatica si convince a tornare indietro e a fare cambio, ma questa prima parte di snorkeling la lascia molto affaticata. Con la sua maschera va molto meglio e anche le nostre esplorazioni sembrano risentire di questo cambiamento, perché troviamo un'altra tartaruga. È ben più grande di quella del giorno prima e sicuramente più anziana. Nuota placidamente, circondata da turisti che le si avvicinano e la toccano sul guscio. Non mi piace molto questo gesto, ma devo ammettere di non aver visto nessuno particolarmente invadente o aggressivo, e in generale si percepisce che queste tartarughe giganti siano abituate ai turisti. Scattiamo foto e video, poi la tartaruga ci lascia per andare a mangiare le alghe nei fondali più profondi.

Continuiamo la nostra esplorazione. C'è una corrente abbastanza forte che spinge verso sud, il che ci mantiene sempre paralleli all'isola, e ci sono le barche ancorate dei tour di snorkeling a cui fare attenzione, per cui stiamo particolarmente vigili anche su quello che succede sopra l'acqua. Il fondale è simile a quello visto il giorno precedente sul lato opposto, fatta eccezione per alcune campane di pietra bucate, probabilmente posate per fare da casa a pesci e altre creature marine. A un certo punto la corrente è troppo forte per tornare via mare, così nuotiamo verso riva e camminiamo lungo il bagnasciuga fino al nostro ombrellone.

Ci tengo a precisare che quando parlo di corrente forte non mi riferisco a niente di insormontabile: il mare non era mosso e la corrente non spingeva verso il largo ma verso sud e un po' verso riva. In ogni caso era pieno di barche ovunque che avrebbero potuto soccorrerci in un'eventuale situazione di pericolo.

Ricongiunti con gli altri raccontiamo subito del nostro incontro e Enri mi propone immediatamente di accompagnarlo alla ricerca di altre tartarughe. Non me lo faccio ripetere due volte e siamo già in acqua. Appena arriviamo al largo mi immergo e davanti a me scorgo subito una tartaruga che nuota. Avviso Enri che si tuffa immediatamente, poi mi fermo: poco più a nord, lungo la linea del reef, ci sono altre due tartarughe — una enorme e più vecchia, una piccola e più giovane. L'euforia è tale che al primo tuffo mi dimentico di compensare e temo di essermi fatto male al timpano. Per fortuna la sensazione passa in fretta, ma questo mi aiuta a essere più consapevole nelle immersioni successive.

Sott'acqua faccio foto e video a volontà: la tartaruga più giovane è molto arzilla, nuota ed esplora curiosando anche nelle zone di acqua più bassa, mentre la più anziana se ne sta a bordo del reef a cibarsi di alghe. Non si tratta di profondità da competizione e posso salire e scendere senza difficoltà, scattando primi piani alle tartarughe e facendo dei bellissimi video a Enrico che nuota con loro. Spostandoci di pochi metri ne troviamo altre: in totale ne avremmo viste circa otto.

Quando torniamo dalle mogli siamo euforici e so già che Elisa mi chiederà di accompagnarla a vedere tutte quelle tartarughe. Ci azzecco — ormai conosco mia moglie — e non faccio in tempo a togliermi le scarpette da roccia che sono di nuovo in acqua.

Il mare nel frattempo si sta ritirando come il giorno precedente e dobbiamo camminare un tratto sempre più lungo prima di poter nuotare. Anche con Elisa siamo fortunati e riesco a filmarla in una danza subacquea con una tartaruga. A un certo punto scorgiamo un esemplare enorme, adagiato sul fondale sul bordo del reef a circa otto metri di profondità. Non giovane e neanche troppo vecchia, ma veramente enorme. Se ne sta lì per i fatti suoi, con l'intenzione evidente di non essere disturbata.

Decido che è il momento di godersi l'animale nel suo habitat naturale senza l'ausilio di mezzi per registrare o fotografare. Chiedo a Elisa di tenermi macchina fotografica e action cam, e mentre lei resta in superficie insieme a un gruppetto di altri italiani, mi preparo rilassando il corpo e ventilando a lungo. Quando mi sento pronto faccio una capovolta e scendo. Devo compensare più volte per arrivare fino alla tartaruga — segno che si trova a circa otto-dieci metri — e quando le arrivo di fianco mi aggrappo con mani e piedi a un enorme corallo duro per non risalire. Lì cerco di rilassare il corpo il più possibile e osservo l'animale nella sua maestosità mentre giace sul fondale. Anche la tartaruga mi osserva con un'espressione di disinteresse, come se volesse dirmi: "Ah, sei qui. Non so chi sei, non mi interessi, basta che non mi importuni." Rilassandomi riesco a rimanere giù per circa un minuto. Sarà un momento da ricordare che porterò a casa con me.

C'è anche una testimonianza video fatta da Elisa in superficie che, come sua abitudine, mi taglia la testa nell'inquadratura.

A proposito delle foto e dei video presi in questa giornata di mare: il corso di apnea mi ha permesso di godere di momenti come questo e di scattare ottime fotografie anche a chi si immergeva con me. Certo, un ricordo di me insieme a una tartaruga, non tagliato, sfocato o mosso, schifo non avrebbe fatto. Che vita grama.

Con questa immersione a tu per tu con la tartaruga mi sono preso comunque una piccola rivincita, a giudicare dall'espressione sconcertata del gruppetto di italiani che mi attendeva all'uscita e che mi ha salutato con un pollice in su.

Riprendiamo a esplorare ma siamo ormai stanchi e quando una tartaruga ci porta nuovamente nel tratto di forte corrente decidiamo di rientrare. Sono abbastanza stanco ma il tempo per me è volato. Mi aspetto che possa essere mezzogiorno, ma in realtà si sono fatte le due e mezza: sono stato in acqua per tre ore.

Abbiamo tutti fame così lasciamo i lettini alla ricerca di un pranzo veloce. Un semplice localino sulla spiaggia fa al caso nostro: le ragazze ordinano pizza margherita, noi due kebab. Non ci sono molte parole per descriverlo — il cibo fa schifo. Pizze surgelate comprate al supermercato e kebab di pollo senza pollo, con insalata che sembra bollita. L'unica cosa degna di nota è la bevanda di Elisa: acqua di cocco con miele, zenzero e zucchero, servita ovviamente all'interno di una noce di cocco.

Terminato il pranzo torniamo di buona lena alla spiaggia per un'ultima immersione. Siamo stanchi e titubanti così prendo al volo l'iniziativa e vado per primo, dicendo a Enrico ed Elisa di seguirmi. Abbiamo solo due maschere e due scarpette da scoglio, così mi avventuro con le infradito. Non è facile, ma almeno raggiunte le acque profonde posso usarle come pinne improvvisate. Entro in acqua mentre gli altri, più veloci, mi raggiungono a piedi.

Si vede che la voglia di bagnarsi a stomaco pieno è poca, perché mi chiedono di andare in avanscoperta da solo e di chiamarli solo se mi fossi imbattuto in una tartaruga. Ma certo, va bene, faccio da esca. Quando raggiungo l'acqua alta immergo la testa, mi guardo attorno, esco e faccio loro il segno del tre con la mano: ci sono tre tartarughe esattamente davanti a me che aspettano solo noi! Enrico ed Elisa, dapprima indecisi, cambiano subito idea e mi seguono a ruota. Io e Enrico, alternandoci la maschera, continuiamo a fare foto e video e a nuotare con questi splendidi animali. Enrico dopo un po' torna indietro lasciando me ed Elisa in acqua.

In breve la corrente cambia direzione e non mi sento più al sicuro così dico a Eli di rientrare. Siamo due che nelle vacanze prendono dei rischi per vivere al massimo, ma si tratta sempre di rischi controllati e questa situazione non è da meno: è tardi, il sole sta per tramontare, siamo stanchi e cotti a puntino, con la pancia piena e l'acqua che inizia a raffreddarsi, con la corrente che cambia nel giro di un attimo. Non ci sono più le condizioni adatte per continuare. Facciamo dietrofront e con un po' di fatica torniamo in spiaggia, percorrendo l’ultimo tratto a piedi tra coralli e stelle marine, dal momento che la marea è ormai rientrata.

Raccogliamo gli zaini e ci dirigiamo verso il centro del paese per prenotare un tavolo per la serata. Andiamo in uno dei locali che la sera prima ci aveva rimbalzato con un "Only drinks" e chiediamo di prenotare. Per le 9 dicono che è troppo tardi e per le 8:30 chiedono un anticipo che suona tanto di mazzetta. Alla fine li convinciamo per le 8 e lasciamo come nome "Marco Polo."

Si è fatto tardi e dobbiamo camminare un bel po' per tornare in albergo, dove dobbiamo anche preparare le borse per l'indomani, quindi ci diamo una mossa. Tornati nella villetta scopriamo che non sono venuti a portare via i piatti sporchi della colazione e nemmeno a fare il refill di birra e acqua potabile dal frigo. Mi preparo in fretta, faccio la borsa per la partenza, copro piatti e tazze sporchi scacciando un geco affamato, e scrivo un messaggio politicamente corretto alla reception chiedendo se per favore possono almeno ricaricare l'acqua potabile per la notte. La speranza è che, una volta entrati nella villa, trovino i piatti sporchi e le birre vuote e portino via tutto. Mi confermano di aver preso in carico la richiesta e, fiducioso si tratti di una semplice svista, non ci penso più.

Ci incontriamo con Enri e Cri e insieme ci dirigiamo verso il ristorante. Elisa rimane indietro a fare video per le stradine buie e intanto racconta quello che vede e quello che abbiamo fatto — sembra una vlogger con la vista a infrarossi. Non ho ben capito perché abbiamo scelto come momento per filmare la sera e come location le vie centrali senza illuminazione: temo che non se ne farà molto di quei video.

Arriviamo puntuali al ristorante, dove non è necessario ricordargli di Marco Polo dal momento che metà tavoli sono vuoti. Ci sediamo e con un po' di difficoltà ordiniamo: sono molto disorganizzati, l'attesa è lunga e altri ospiti arrivati dopo di noi vengono serviti prima. Per fortuna siamo intrattenuti dalla musica live di una band che alterna brani famosi a musica tradizionale. Mi colpisce molto un anziano nativo americano dai capelli lunghi e grigi che, dopo aver suonato il flauto, intona una canzone tradizionale. Eli rimane così colpita da andare in prima fila a fianco del palco per filmare la scena da vicino.

Il posto è turistico e molto frequentato e anche stavolta la cena sulla spiaggia non soddisfa le aspettative. Abbiamo la possibilità di riflettere su un problema abbastanza comune: la grande attenzione e l'infinita gentilezza del personale non compensano la forte disorganizzazione. Durante le prime tappe del viaggio siamo spesso stati gli unici clienti o comunque in gruppi molto ridotti: le portate venivano preparate una per una, con calma e attenzione, e questo faceva percepire di essere coccolati. Quando abbiamo incontrato i primi contesti con molte persone da gestire sono arrivati i primi problemi.

È così che stasera il mio piatto principale arriva prima degli antipasti e quando faccio notare che probabilmente si tratta dell'ordine di un altro tavolo, il cameriere mi consiglia comunque di prenderlo per non rischiare di aspettare ancora più a lungo. Sono molto confuso da tale logica, ma ho fame e acconsento. Per fortuna a breve arrivano anche gli antipasti insieme alle altre portate e col tavolo pieno possiamo finalmente banchettare, salvo l'ennesima sorpresa per Elisa: i suoi spiedini di pollo Satay sono ricoperti di salsa piccantissima. Per permetterle di mangiare qualcosa le propongo di fare a cambio, ma il mio piatto essendo arrivato prima degli altri è già praticamente finito. Accetta lo scambio lo stesso e io cerco di mangiare il suo pollo piccantissimo rimuovendo la salsa come meglio riesco: il risultato è un piatto di riso in bianco con pollo scondito e sporadiche punte di piccante.

Terminata la cena andiamo subito a pagare prima che facciano altri danni. Mi fermo in bagno e inizialmente rimango stupito dall'affresco sul soffitto che ricorda molto i cieli di Van Gogh, ma poi inorridisco quando realizzo che si tratta di muffa. Fuggo via.

La serata prosegue passeggiando per le vie dell'isola, soffermandosi davanti ai locali con musica dal vivo. Assistiamo alla live di una band il cui frontman è un indonesiano mingherlino in canotta e jeans neri, che chiude la serata cantando a squarciagola \textit{Highway to Hell}, attirando l'attenzione di tutti i passanti. Incrociamo per strada un gruppo con le t-shirt tutte uguali con la scritta "Pub Crawl Gili Trawangan", che riconferma il target del turista medio e aiuta ad allontanare la malinconia dell'ultima sera.

Ci siamo trovati davvero bene su quest'isola, soprattutto quando siamo riusciti a viverla a modo nostro evitando il turismo di massa dei tour, delle discoteche e dell'alcolismo da pub crawl. Viene il dubbio che Gili Air e Gili Meno condividessero la stessa bellezza e le stesse attrattive senza avere quell'anima festaiola che qui abbiamo costantemente evitato.

Siamo tutti stanchi e arriva presto il momento di salutarci. Ci abbracciamo e ci auguriamo un'ottima continuazione di vacanza, dandoci appuntamento a Torino.

Rientrati nella villa purtroppo ci accoglie un'amara sorpresa: non solo non hanno tolto piatti sporchi e bottiglie vuote, ma non ci hanno nemmeno portato l'acqua potabile per la notte. Mi accingo ad andare alla reception a farmi sentire, ma non trovo nessuno — solo il distributore automatico dell'acqua, che almeno ci permette di andare a dormire con le scorte. A letto per l'ultima volta su quest'isola paradisiaca.


\pagestyle{empty} % disable numbers

%coppia
\landscapeLargeMargin{images/day10/PXL_20240818_014501139.MP.jpg}
%

%coppia
\landscapeLargeMargin{images/day10/PXL_20240818_024359146.jpg}
\landscapeLargeMargin{images/day10/PXL_20240818_030015328.jpg}
%

%coppia
\landscapeLargeMargin{images/day10/PXL_20240818_024359146.jpg}
\landscapeLargeMargin{images/day10/PXL_20240818_030015328.jpg}
%

%%% TODO: faccia ritagliata della tartaruga
%coppia
\splitImageLeft{images/Kodak/Screenshot 2026-05-19 at 22.36.05.png}
\splitImageRight{images/Kodak/Screenshot 2026-05-19 at 22.36.05.png}
%

%coppia
\landscapeThinMarginCentered{images/Kodak/Screenshot 2026-05-19 at 21.46.14.png}
\landscapeThinMarginCentered{images/Kodak/Screenshot 2026-05-19 at 21.40.29.png}
%

%coppia
\twoImagesThinMargins{images/day10/Screenshot 2026-05-18 at 18.46.04.png}
                      {images/day10/PXL_20240818_screen001.png}
\landscapeThinMarginCentered{images/day10/Screenshot 2026-05-18 at 18.47.49.png}
%

%coppia
\splitImageLeft{images/Kodak/Screenshot 2026-05-19 at 21.36.15.png}
\splitImageRight{images/Kodak/Screenshot 2026-05-19 at 21.36.15.png}
%

%coppia
\landscapeThinMarginCentered{images/day10/Screenshot 2026-05-18 at 18.44.36.png}
\twoImagesThinMargins{images/day10/Screenshot 2026-05-18 at 18.49.06.png}
                      {images/day10/Screenshot 2026-05-18 at 18.49.35.png}
%

%coppia
\twoImagesThinMargins{images/Kodak/Screenshot 2026-05-19 at 21.43.01.png}
                      {images/Kodak/Screenshot 2026-05-19 at 21.39.02.png}
\landscapeThinMarginCentered{images/day10/Screenshot 2026-05-18 at 18.50.05.png}
%

%coppia
\threeImagesThinMargins{images/day10/PXL_20240818_075409711.jpg}
                       {images/day10/PXL_20240818_072644892.MP.jpg}
                       {images/day10/PXL_20240818_072646504.MP.jpg}
\landscapeLargeMargin{images/day10/PXL_20240818_101947796.jpg}

%single
\twoImagesLargeMargins{images/day10/PXL_20240818_124623817.MP.jpg}
                      {images/day10/PXL_20240818_124710154.MP.jpg}
%

\pagestyle{fancy} % enable numbers 




